% 依据《全国大学生数学建模竞赛人工智能工具使用规定》（2026 年试行）第 3 条：
% 「参赛队应在论文**参考文献之前**设置"AI 工具使用声明"，二者择一。」
% 本文属于「使用了 AI 工具」的情形，采用第（2）种声明。
% 详细使用情况另见支撑材料中的《AI 工具使用详情.pdf》（第 4 条要求）。

\section*{AI 工具使用声明}
\addcontentsline{toc}{section}{AI 工具使用声明}

本参赛队在竞赛过程中使用了 AI 工具，主要用于\textbf{代码调试与数值结果核验、
文献检索与出处核实、以及语言润色}，详细使用情况见支撑材料
《AI 工具使用详情.pdf》。

\vspace{2pt}
\noindent\textbf{核心建模与分析的自主性声明}：

\begin{enumerate}[leftmargin=2.2em, itemsep=2pt]
  \item 本文的\textbf{物理模型、控制方程、边界条件、数值格式与判据口径}
        均由参赛队自主确定；AI 工具不参与模型选择与建模决策。
  \item 本文\textbf{全部数值结果}由参赛队编写的求解程序计算产生。
        AI 工具用于代码调试与结果核验，其输出\textbf{逐项经过人工审查}：
        凡由 AI 辅助发现的代码缺陷，均要求提供\textbf{可复现的最小反例}后
        才予采纳（例如"误差随步长收敛阶被 $\log 0$ 污染而虚高"、
        "台账合并误删历史记录"两处，都是先复现再修）。
  \item 本文\textbf{所有引用的数据与结论均可溯源}：
        题面与附件数据标【题面/附件】，文献结论标明出处与 DOI，
        本文自算量显式标注"本文计算得到"。
        \textbf{凡无法核实的外部数字一律不采用}（见 \S\ref{sec:refnote}）。
  \item AI 工具生成的文字\textbf{均经人工改写与核验}，
        并由参赛队对全文的数值一致性做了\textbf{自动校验}
        （逐条比对论文中的关键数值与结果文件，共 25 项）。
\end{enumerate}
